\newpage
%===============================================================================================
% Event's programs, information or other content
%===============================================================================================

% It's important to include "\hoffset=0pt"
% before including PDF's or other additional files, as the "\setlength{\hoffset}{-0.5in}"
% in the .cls moves all the included files to the left
{\hoffset=0pt
\includepdf[fitpaper=true,pages=-,width=\paperwidth]{imagenes/eventos/sipa xiii programa.pdf}
}

{\hoffset=0pt
\includepdf[fitpaper=true,pages=-,width=\paperwidth]{imagenes/eventos/sipa xiii curso 1.pdf}
}

{\hoffset=0pt
\includepdf[fitpaper=true,pages=-,width=\paperwidth]{imagenes/eventos/sipa xiii infografias.pdf}
}

{\hoffset=0pt
	\includepdf[fitpaper=true,pages=-,width=\paperwidth]{imagenes/eventos/sipa xiii conversatorio.pdf}
}

%======================
%	Geometry
%======================

\bookbodygeometry

% Set roman page numbering, custom page format and page number
\pagenumbering{roman}
\pagestyle{preindex}
\setcounter{page}{9}

%===============================================================================================
% Event's ntroduction
%===============================================================================================

\subseccion{prólogo}

El evento de la SEMANA INTERNACIONAL EN PRODUCCIÓN AGROPECUARIA se realiza año con año fortaleciendo las actividades de divulgación científica, tiene la finalidad de fomentar la difusión sistemática de los trabajos realizados por investigadores institucionales, nacionales, así como por extranjeros. En cada celebración del evento, además de dar a conocer los trabajos de investigación científica de reconocidos investigadores, se promueve la vocación científica con la participación activa de nuestros estudiantes, quienes en conjunto con sus asesores, dan a conocer avances de sus proyectos de investigación o plantean nuevas propuestas con el rigor científico y el aporte suficiente para ser presentados, los aportes están claramente relacionados con temas de relevancia científica que contribuyen significativamente a la solución de problemas y retos de interés general, orientados a mejorar la competitividad y la productividad de los sectores económicos del país, ayudando así a incrementar la calidad de vida de la población y del medio ambiente que nos rodea. Pese a la situación atípica que estamos viviendo por pandemia por COVID, los trabajos presentados en esta VIII SEMANA INTERNACIONAL EN PRODUCCIÓN AGROPECUARIA, son el reflejo de que se mantiene acercamiento de investigadores y estudiantes con los diversos sectores sociales y productivos. Además de que se sostiene el compromiso institucional y del programa de posgrado para con la sociedad y el medio ambiente, para continuar con propuestas y desarrollo de proyectos de investigación para la atención de las principales problemáticas científicas ambientales y sociales, de ámbito regional y local que enfrenta el país y sus referentes en el mundo. La celebración de la SEMANA INTERNACIONAL EN PRODUCCIÓN AGROPECUARIA es y seguirá siendo un espacio para el reconocimiento público de la investigación y la innovación en la incidencia en el bienestar social y la sustentabilidad. La realización anual de dicho evento contribuye en el fortalecimiento de la cultura científica y tecnológica en la sociedad, por lo que se considera una actividad de retribución social que forma parte integral del programa de Postgrado en Ciencias en Producción Agropecuaria.

\raggedleft
Dra. Ma. Guadalupe Calderón Leyva\\[1mm]
Jefe del Programa de Posgrado en Ciencias en Producción Agropecuaria\\[2mm]
Prólogo de la VIII SIPA

\raggedright
\newpage

%===============================================================================================
% Event's directory
%===============================================================================================

\subseccion{directorio}
\vspace*{\fill}
\begin{center}
	\fontsize{13}{16}\selectfont
	
	\textbf{Rector}
	
	X\\[1.8cm]
	
	\textbf{Director Regional}
	
	M. C. Edgardo Cervantes Álvarez\\[1.8cm]
	
	\textbf{Subdirector de Postgrado}
	
	X\\[1.8cm]
	
	\textbf{Jefe del Departamento de Postgrado}
	
	Dra. Dalia Ivette Carillo Moreno\\[1.8cm]
	
	\textbf{Jefe del Programa de Maestría en Ciencias en Producción Agropecuarias}
	
	Dra. Viridiana Contreras Villarreal\\[1.8cm]
	
	\textbf{Jefe del Programa de Doctorado en Ciencias en Producción Agropecuarias}
	
	Dra. Ma. Guadalupe Calderón Leyva\\[1.8cm]
	
\end{center}
\vspace*{\fill}
\newpage

%===============================================================================================
% Organizing committees
%===============================================================================================

\begin{multicols}{2}
[\subseccion{Comités}]
\centering
\subtitulos{Presidente organizador}
X

\subtitulos{Comité organizador}
X

\subtitulos{Comité técnico}

\subtitulos{Comité organizador de memorias}
Dr. Francisco Gerardo Véliz Deras

Dra. Dalia Ivette Carrillo Moreno

Dra. Ma. Guadalupe Calderón Leyva

Dra. Jessica Flores Salas

Dra. Viridiana Contreras Villarreal

Dr. Cayetano Navarrete-Molina

\subtitulos{Comié evaluador de carteles}
X

\subtitulos{Comité de registro}
X

\subtitulos{Moderadores}
X

\subtitulos{Comité de apoyo logístico}
Sra. Aurelia Nájera Cruz

\subtitulos{Colaboradores}
Alumnos de Maestría en Ciencias en Producción Agropecuaria

Alumnos de Doctorado en Ciencias en Producción Agropecuaria

\end{multicols}
